\documentclass[11pt,a4paper]{article}
\usepackage[margin=21mm,headheight=15pt]{geometry}
\usepackage[T1]{fontenc}
\usepackage{lmodern,amsmath,amssymb,booktabs,tabularx,enumitem,xcolor,fancyhdr}
\usepackage[hidelinks]{hyperref}
\definecolor{accent}{RGB}{0,114,178}
\setlength{\parindent}{0pt}
\setlength{\parskip}{6pt}
\setlist{nosep,leftmargin=*}
\pagestyle{fancy}
\fancyhf{}
\fancyhead[L]{\small Econometrics 25117 | Instructor guide}
\fancyhead[R]{\small 2026--2027}
\fancyfoot[L]{\small Private teaching notes}
\fancyfoot[R]{\thepage}
\newcommand{\heading}[1]{\section*{\color{accent}#1}}
\newcommand{\slideheading}[1]{\subsection*{\color{accent}#1}}
\newcommand{\cue}[2]{\par\noindent\textbf{#1}\quad #2\par}
\newcommand{\slideguide}[3]{\slideheading{Slide #1 | #2}\cue{Guide}{#3}}
\setlength{\emergencystretch}{1em}
\newcommand{\E}{\mathbb{E}}
\newcommand{\Var}{\operatorname{Var}}
\newcommand{\Cov}{\operatorname{Cov}}
\begin{document}
\begin{center}{\Large\bfseries Inference about means}\end{center}
\textbf{Slide route:} Lecture2v2.tex: hypothesis testing through California Test Score Data Set. Stock and Watson Chapter 3.
\heading{Session 3 | Evidence against a claim}
\textbf{Outcomes:} State a null and alternative, compute a test statistic and confidence interval, distinguish statistical from economic significance, and compare two independent groups.
\cue{Opening}{A sample mean is 104; a policy report claims the population mean is 100. Is four units a large discrepancy? We need the scale of sampling uncertainty.}
\heading{100-minute route}
\begin{tabularx}{\textwidth}{@{}rX@{}}\toprule
0--10 & Recall sampling distributions and SE.\\
10--30 & Hypothesis Testing, Key Concepts, Computing the P-Value.\\
30--45 & Sample variance, standard errors, t-statistic; worked test.\\
45--55 & Type I/II errors and short pause.\\
55--70 & Confidence intervals and test/interval correspondence.\\
70--90 & Mean Difference Testing, randomization, California application.\\
90--100 & Exit exercise and preview of a regression slope.\\\bottomrule
\end{tabularx}
\heading{Slide-by-slide script}
\slideguide{1}{Title}{Open with the question: how do we decide whether a population claim is compatible with sample evidence?}
\slideguide{2}{Roadmap}{Tell students that today moves from a null hypothesis to a test, an interval, and finally a comparison of groups.}
\slideguide{3}{Last lesson: probability}{Retrieve one fact: $\E(Y\mid X=1)$ is a conditional mean, not the unconditional mean. Use the sunny-day example for a 30-second calculation.}
\slideguide{4}{Last lesson: sampling}{Ask why $\Var(\bar Y)=\sigma^2/n$ matters for inference. Students should answer: it determines the scale of sampling noise.}
\slideguide{5}{What is hypothesis testing?}{State $H_0$ and $H_1$ before looking at data. Example: test $H_0:\mu=100$ against $H_1:\mu\neq100$.}
\slideguide{6}{What is hypothesis testing?}{Emphasize that the null is a reference model, not automatically the truth. Ask what a two-sided alternative allows.}
\slideguide{7}{Why use hypothesis testing?}{Connect the procedure to a decision under uncertainty. Example: a drug trial where a sample improvement may be sampling variation.}
\slideguide{8}{Key concepts}{Define test statistic, reference distribution, and p-value in that order. Ask which object is random under the null: the statistic, not the fixed parameter.}
\slideguide{9}{Test statistic setup}{Put the observed mean and hypothesized mean on a number line. The distance must be scaled by its standard error.}
\slideguide{10}{Computing the p-value}{Use the two-sided rule: count probability in both tails at least as far from zero as the observed statistic. Warn against reporting $P(H_0\mid\text{data})$.}
\slideguide{11}{Sample variance}{Calculate $s^2=(n-1)^{-1}\sum_i(y_i-\bar y)^2$ and explain why the denominator is $n-1$.}
\slideguide{12}{Standard error}{For $n=100$, $s=20$, calculate $SE=2$ on the board. Ask what happens to $SE$ if $n$ becomes 400.}
\slideguide{13}{Type I and Type II errors}{Use a pregnancy-test analogy only to label false positive and false negative. Then return to econometrics: Type I rejects a true null; Type II fails to reject a false null.}
\slideguide{14}{Errors in testing}{Explain the trade-off between a smaller significance level and power. Do not call failure to reject acceptance of the null.}
\slideguide{15}{Confidence intervals}{Construct the interval as estimate plus and minus critical value times standard error. Ask what repeated-sampling coverage means.}
\slideguide{16}{What are CIs?}{Use the interval to show that a two-sided 5 percent test and a 95 percent interval answer matching questions under the same reference distribution.}
\slideguide{17}{Computing CIs}{Work $104\pm1.96(2)=[100.08,107.92]$. Ask whether the interval contains 100 and what that implies for the corresponding test.}
\slideguide{18}{Mean difference testing}{Define $\widehat\Delta=\bar Y_1-\bar Y_0$ and keep the group labels visible. Example: $70-65=5$.}
\slideguide{19}{Mean difference testing}{Compute $SE(\widehat\Delta)=\sqrt{s_1^2/n_1+s_0^2/n_0}$. Ask why independent groups permit variances to add.}
\slideguide{20}{Randomized control trials}{Explain why random assignment makes a difference interpretable as a causal effect under the design. It does not guarantee equal groups in one finite sample.}
\slideguide{21}{California data: setup}{Name the outcome, treatment or regressor, unit, and sample before interpreting any coefficient.}
\slideguide{22}{California data: descriptive comparison}{Ask students to compare average test scores by class size category without calling the gap causal.}
\slideguide{23}{California data: uncertainty}{Identify the estimate, standard error, and confidence interval in the displayed output. Have students point to each rather than reading the table aloud.}
\slideguide{24}{California data: test}{Choose one null, write it symbolically, and calculate the corresponding statistic. Keep the null visible while interpreting.}
\slideguide{25}{California data: interpretation}{Separate statistical significance from economic importance. Ask whether the interval rules out effects that would matter for policy.}
\slideguide{26}{California data: limitations}{List selection, omitted variables, and measurement concerns. Correct arithmetic does not by itself create identification.}
\slideguide{27}{California data: conclusion}{Have students give a complete sentence with estimate, units, uncertainty, and causal qualification.}
\slideguide{28}{Review}{Use the synthetic two-group example from this guide as a one-minute retrieval exercise.}
\slideguide{29}{Material}{Tell students which chapter pages and problem set questions to use for practice; do not spend class time reading the bibliography.}

\heading{Board example | A mean}
\cue{Use with slides}{Lecture2v2 slides 5--11: What is Hypothesis Testing? through the t-statistic. Start the numerical mean test on slide 8, compute the standard error on slide 10, and finish the test statistic on slide 11.}
Illustrative sample: $n=100$, $\bar y=104$, $s=20$. Test $H_0:\mu=100$ against $H_1:\mu\ne100$.
\[
SE(\bar Y)=20/\sqrt{100}=2,\quad t=(104-100)/2=2.
\]
Using a large-sample normal approximation, the two-sided $p$-value is about $.0455$ and a 95\% interval is
\[
104\pm1.96(2)=[100.08,107.92].
\]
Reject at 5\% but not at 1\%. If observations are i.i.d. normal and variance is estimated, an exact $t_{99}$ reference gives a slightly larger $p$-value (about $.048$). State which reference distribution is being used.
\cue{What to say}{A $p$-value is the probability, under the null and assumptions, of a statistic at least as extreme as the observed one. It is not the probability that the null is true.}
\cue{If short}{Keep the mean example and two-group example. Use one California output table deeply; assign repetitive descriptive tables for review.}

\newpage
\heading{Side example | Comparing groups}
\cue{Use with slides}{Lecture2v2 slides 15--19: Confidence Intervals, Mean Difference Testing, and Causal Inference in Randomized Control Trials. Use the two-group calculation on slides 17--18, then discuss the random-assignment caveat on slide 19.}
Two independently sampled groups have means $70$ and $65$, SDs $10$ and $12$, and sample sizes $50$ each:
\[
\widehat\Delta=5,\quad SE(\widehat\Delta)
=\sqrt{10^2/50+12^2/50}=\sqrt{4.88}\simeq2.209.
\]
\[
t\simeq2.264,\quad 95\%\text{ large-sample CI: }
5\pm1.96(2.209)=[.67,9.33].
\]
Ask students what changes if these are the same individuals measured twice. The independent-groups formula is then inappropriate; the covariance between measurements matters.

\cue{Causal interpretation}{If groups are assigned randomly and the outcome comparison follows that assignment, a mean difference can estimate a treatment effect under the design's assumptions. If students choose their own school, the same arithmetic does not remove selection.}
\cue{Type I/II prompt}{A Type I error rejects a true null. A Type II error fails to reject a false null. Holding sample size and effect fixed, demanding a smaller significance level generally reduces power.}
\cue{Confidence-interval language}{The method covers the fixed population parameter in 95\% of repeated samples under its assumptions. Avoid saying that 95\% of observations lie in this interval.}
\cue{Economic significance}{An enormous sample can make a tiny effect statistically significant. Ask whether the interval excludes effects small enough to be irrelevant to the decision.}
\heading{Exit exercise and answers}
If $n=400$ with the same SD and the same observed gap of four units, $SE=1$ and $t=4$. The standard error halves, not quarters.

If a 95\% interval includes zero, fail to reject a two-sided zero-effect null at the corresponding 5\% level. This is not proof of zero effect.

\cue{Next preparation}{Read Chapter 4. Bring a scatterplot of test scores against class size to the next class and ask how to summarize its slope.}
\textbf{Source:} Lecture2v2.tex. The numerical data above are synthetic teaching examples.

\end{document}
